\section{Installing and running}
\label{sec:installation}

\subsection{Requirements}

\begin{center}
\begin{tabular}{ll}
\toprule
Component & Version and role \\
\midrule
Python & $\geq 3.10$ \\
numpy, scipy & required; the analysis layer needs only these \\
pyscfad & $\geq 0.3.3$; provides the JAX-based \texttt{gto} and \texttt{scf} \\
pyscf & $\geq 2.5$; reference comparisons and integral drivers \\
jax & $\geq 0.4$; the differentiable array layer \\
matplotlib & only for the figure sources \\
\bottomrule
\end{tabular}
\end{center}

\texttt{pyscfad} pins a specific JAX version, so install the package
into a dedicated environment rather than into a shared one.

\begin{lstlisting}[language=bash]
conda create -n diffmrc python=3.11
conda activate diffmrc
pip install -e ".[core,dev]"     # kernel, tests and lint tooling
\end{lstlisting}

For work on the analysis layer alone, \texttt{pip install -e ".[dev]"}
is sufficient: that layer is plain NumPy and imports nothing from the
numerical kernel.

\subsection{Runtime settings that are not optional}

Two settings must be established before the first JAX import.

\begin{lstlisting}[language=Python]
import os
os.environ['JAX_PLATFORMS'] = 'cpu'     # before importing jax
for v in ('OMP_NUM_THREADS', 'MKL_NUM_THREADS',
          'OPENBLAS_NUM_THREADS', 'NUMEXPR_NUM_THREADS',
          'VECLIB_MAXIMUM_THREADS'):
    os.environ[v] = '8'

import jax
jax.config.update('jax_enable_x64', True)
\end{lstlisting}

\texttt{jax\_enable\_x64} must be set: the package is validated in
double precision, and single-precision runs will produce silently
wrong gradients.  The CPU platform pin is deliberate.  The stack has
been developed and verified on CPU only; on a machine with a GPU the
default JAX platform would select CUDA, where this build of pyscfad
raises a solver error that has no bearing on the chemistry.  Pin the
platform for now; GPU support is not claimed.

The thread variables keep the BLAS layer inside the thread budget the
driver assumes.  They are not a performance suggestion: without them a
single calculation can oversubscribe a shared node.

\subsection{Checking the installation}

\begin{lstlisting}[language=bash]
pytest tests/ -v        # unit and integration tests
lint-imports            # verifies the layer boundaries
\end{lstlisting}

The layer contracts are worth a word even at install time.  The package
is organized in six layers (kernel, analysis, policy, driver, report,
command line), and the import direction between them is enforced as a
test.  If \texttt{lint-imports} is unavailable in your environment,
\texttt{tests/test\_layer\_contracts.py} performs the same check.  A
run that passes the suite but fails the contracts is a run the
developers would want to see.

Two environment notes that cost us time, recorded so that they need
not cost yours:

\begin{itemize}
  \item pyscfad's molecule object cannot be traced inside
    \texttt{jax.jit}; the driver layer constructs plain-array callables
    for anything that is jitted.  If you write your own jitted wrapper
    around a factory from \texttt{driver.nevpt2}, use the
    \texttt{theta}-only factories, which are built for that context.
  \item The first call into a heavy node of the chain compiles an XLA
    program and can take minutes; subsequent calls are fast.  A slow
    first evaluation is expected, not a hang.
\end{itemize}
